% =============================================================================
%  Institutional Profiles Under a Common Prudential Framework:
%  Evidence from Brazilian Credit Cooperatives and Banks
%  ICA CCR 2026 full submission, Best Young and Emerging Scholar Paper
%  Word limit: 8,000 including tables, figures and references.
%
%  Tables come from tables/*.tex, generated by make_tables.py from results/*.txt.
%  Figures come from figures/. Do not transcribe numbers by hand.
%  Tables t1, t2_main_ci, t3, t4, t5, t6, t8, t9, t11, t12 and figures fig10, fig15,
%  fig16 and fig17 live in supplement.tex. Figures come from make_figures.py.
% =============================================================================
\documentclass[11pt,a4paper]{article}
\usepackage[utf8]{inputenc}
\usepackage[T1]{fontenc}
\usepackage{lmodern}
\usepackage[margin=1in]{geometry}
\usepackage{booktabs}
\usepackage{graphicx}
\usepackage{amsmath}
\usepackage{threeparttable}
\usepackage{microtype}
\PassOptionsToPackage{hyphens}{url}
\usepackage[hidelinks]{hyperref}
\usepackage[numbers,square,sort&compress]{natbib}
\usepackage{setspace}
\onehalfspacing

\title{\textbf{Institutional Profiles Under a Common Prudential Framework:}\\[0.3em]
\large Evidence from Brazilian Credit Cooperatives and Banks}
\author{Arthur Gomes Nery$^{a}$, Thiago de Oliveira Victorino$^{a,b}$ and
Rodrigo Lima Rangel$^{a}$\\[0.6em]
\small $^{a}$Organization of Brazilian Cooperatives (OCB System)\\
\small $^{b}$Department of Administration, University of Brasilia, Brazil}
\date{}

\begin{document}
\maketitle
\vspace{-2em}
\begin{center}
\textit{Award category:} Best Young and Emerging Scholar Paper
\end{center}

\begin{abstract}
\noindent In Brazil, a credit cooperative that grows past a size threshold computes
capital, leverage and risk under the same prudential methodology as a commercial bank. We
use this setting to ask whether the cooperative form retains distinct balance-sheet
characteristics. Using a quarterly panel of every
full-methodology institution from 2017 to 2024, 130 singular credit cooperatives and 144
banks, we relate ten measures of capital, performance, stability and balance-sheet
structure to a cooperative dummy. The estimates control for size and time and are
repeated on a matched sample and across five comparison groups. Three results emerge. First, cooperatives of comparable size allocate more of
their assets to credit than banks nationally, operate at lower cost, show lower return
volatility and earn a higher return on assets after tax, a gap that the tax exemption of
cooperative acts explains in part. Second, capital differs modestly in level and
markedly in dispersion. The mean gap comes largely from small banks that barely
intermediate and hold capital ratios above fifty percent, and it roughly halves against
deposit-funded banks. Among thinly capitalised institutions, cooperatives hold more
capital than banks and provision no less. Third, lower provisioning and narrower margins,
two differences often reported, disappear once the comparison group is restricted to
banks. A benchmark calibrated on the average institution may therefore describe the
locally lending cooperative poorly. Whether that bears on the territorial role of
cooperative lenders is the claim these data cannot yet settle.

\medskip
\noindent\textbf{Keywords:} credit cooperatives; prudential regulation; Basel;
proportionality; Brazil.\\
\textbf{JEL:} G21, G28, P13.
\end{abstract}

\newpage

% =============================================================================
\section{Background}
% =============================================================================

Prudential regulation has converged on a common template. Basel III supplies the
architecture for capital, leverage and risk measurement, and national supervisors adapt
it to local conditions \citep{bcbs2011}. Brazil sorts authorised institutions into five
segments for the proportional application of prudential rules \citep{cmn4553};
Figure~\ref{fig:regulatory} shows how cooperatives and other institutions divide between
the two methodologies. Segments 1 to 4 apply the full methodology, while Segment~5
comprises institutions with total exposure below 0.1~percent of GDP that opt for a
simplified method of computing minimum capital, under Resolution 4.194/2013 until February 2018 and
Resolution 4.606/2017 thereafter \citep{cmn4606}. Commercial, multiple,
investment and exchange banks are excluded from S5 at any size.\footnote{The segment
label is revised semiannually, so five full-methodology
cooperative-quarters still carry the S5 label in transition. The sample follows the
reported methodology.} Institutions outside S5 apply the \textit{Patrim\^onio de
Refer\^encia} capital definition, minimum requirements and risk-weighted-asset
computation of Resolutions 4.192 and 4.193 of 2013 and, from 3 January 2022, of
Resolutions 4.955 and 4.958 \citep{cmn4955, cmn4958}. Loan
classification and minimum provisioning follow Resolution 2.682/1999 throughout the
window, since its expected-loss successor took effect only in 2025 \citep{cmn2682,
cmn4966}.

A credit cooperative that grows past the threshold, or declines the simplified option,
therefore computes capital under the same definition, minimum ratios, standardised risk
weights and provisioning rules as a commercial bank. A few requirements differ by
segment: the liquidity ratios and the systemic buffer apply to S1 only, and the Basel
leverage ratio to S1 and S2 \citep{cmn4401, cmn4616, cmn4958, cmn4615}. Those segments
contain banks only, and Section~\ref{sec:robustness} shows that excluding them changes
little.

\begin{figure}[htbp]
\centering
\includegraphics[width=\textwidth]{figures/fig5_regulatory.png}
\caption{Institutions by prudential methodology, 2017 to 2024}
\label{fig:regulatory}
\par\smallskip{\footnotesize\raggedright Panel (c) is the share under the full methodology. The 2023 step in panel (b) is the entry of payment-institution conglomerates.\par}
\end{figure}

Credit cooperatives are member-owned deposit-taking institutions that lend mostly to
their members. They are organised under Complementary Law 130/2009 and Resolution
4.434/2015 \citep{lc130, cmn4434}. Cooperatives held about six percent of system assets at the
end of 2024 and 6.3 percent a year later, and their share has risen in every year since
2020 \citep{bcb2025panorama, bcb2026panorama}. Much of the sector belongs to integrated
systems in which central cooperatives supply liquidity facilities, consolidated risk
management and a single supervisory interface to their affiliates.

Cooperatives are the part of the regulated system that lends locally and has grown
fastest. If the prudential benchmark reflects an average institution that does not lend,
it measures cooperatives against a standard drawn from a different kind of institution.
Whether it does is the question we ask, and the answer bears on institutional resilience
and on the territorial role of cooperative lenders.

\subsection{Theoretical framework and expectations}
\label{sec:theory}

Prudential capital regulation exists to internalise the externalities that a financial
institution's failure imposes on depositors, on the payment system and on the safety
net \citep{bcbs2011}. It also offsets the moral hazard that a deposit guarantee
introduces into the institution's choice of leverage \citep{diamonddybvig1983}. In
Brazil the guarantee is mutual and split by form: bank deposits are covered by the FGC
and cooperative deposits by the FGCoop, a private fund created in 2013, capitalised by
the cooperatives at a flat rate that ignores each member's risk, with no central bank
lender of last resort behind it \citep{coelho2019}. One reading of the prudential rationale is a
representation argument, in which the supervisor acts on behalf of small, dispersed
depositors who cannot monitor a leveraged intermediary. In a credit cooperative the
depositor is also a member with a residual claim and a governance right, and in the
integrated systems a central body supplies consolidated monitoring. The monitoring
problem is therefore likely to be partly internalised. The systemic externality of
failure and the moral hazard of the guarantee remain, and under a flat-rate mutual fund
a thinly capitalised cooperative shifts part of its risk onto the rest of the sector. A
capital template written for one form may thus bind, when applied to another, on a
margin that does not carry the risk it was designed to price.

Brazil's segmentation sorts institutions by size and international activity
\citep{cmn4553}, and the international surveys of proportionality record size,
complexity and business model as the criteria in use, none based on organisational form
\citep{bcbs2019, castrocarvalho2017}. Where supervisors adapt the framework to
cooperatives, the adaptation concerns member shares and network arrangements
\citep{coelho2019}. Responsive regulation \citep{ayres1992,
baldwin2008} treats the structure of the regulated industry and the motivations of its
firms as inputs to how a regime engages them. Organisational form and network membership then become
candidate calibration margins, and the profile of each form decides whether they point to
lighter treatment or closer scrutiny. This paper provides that profile.

Ownership offers a reason to expect the two forms to differ under a common methodology. A
cooperative cannot raise outside voting equity as needed, so its capital is an endowment
of retained earnings accumulated across member generations \citep{fonteyne2007}. Basel
III admits the member shares of mutuals and cooperatives as common equity under criteria
adapted to their legal form \citep{bcbs2011}. However, the shares remain redeemable
claims of members and cannot be placed with outside investors. The member's claim carries
the five property-rights problems that \citet{cook1995} sets out: free-rider, horizon,
portfolio, control and influence costs, which together weaken the incentive to build
capital beyond what member service requires, since the member's claim is a claim on the
use of the institution.

\subsection{Comparison groups in the cooperative banking literature}

Besides the proportionality work above, two strands of the literature reviewed by
\citet{mckillop2020} bear on this paper. The first concerns capital, profitability and
stability. \citet{hesse2007} find in a cross-country panel that cooperative banks are
more stable than commercial banks. They attribute the result to the lower volatility of
cooperative returns, which more than offsets lower profitability and capitalisation.
\citet{iannotta2007} find that ownership structure predicts risk and performance across
European banks, and \citet{becchetti2016} and
\citet{ferri2014} show that it shapes lending. \citet{groeneveld2020} weighs the
adaptations of European cooperative banks against the traits they retain and finds
their distinctive character intact. \citet{cuevas2006} set out the tension
between member orientation and prudential templates written for investor-owned banks. The
second is Brazilian. \citet{bressan2011} apply the PEARLS system to Brazilian credit
cooperatives and model their insolvency risk, and \citet{carvalho2015} analyse the exit
and failure of Brazilian credit unions. None of these studies compares cooperatives with
banks under a common prudential methodology, which is the setting we use here.

A further question is asked less often: which institutions the comparison group contains.
Regulatory populations are defined by legal category and can include institutions with
little economic resemblance to a bank, as Section~\ref{sec:peers} shows for Brazil.

\subsection{Study design and contribution}

We compare cooperatives with banks across a ladder of comparison groups defined on the
regulator's own categories, and separate sector-wide differences from differences within
regional markets. The window opens in 2017Q1, the first quarter with segment reporting
under Resolution 4.553, so the paper cannot speak to the crisis-period lending that
\citet{becchetti2016} and \citet{ferri2014} study.

% =============================================================================
\section{Aims}
% =============================================================================

We ask whether organisational form leaves a measurable imprint on the balance sheet once
the prudential methodology is common to both forms, and how far the answer depends on
choices the analyst makes. Section~\ref{sec:theory} leads us to expect cooperatives to
hold less regulatory capital than banks of comparable size, in a narrower band than a
population that can issue equity on demand. Because the member's claim is a claim on use,
we also expect them to lend a higher share of assets. On profitability and efficiency the
framework does not commit to a sign. Surplus returned through pricing would depress the
return on assets, while the cost discipline of a focused retail model would raise it, and
the tax exemption of cooperative acts raises the after-tax return mechanically, so we
examine the return before tax as well. The framework does lead us to expect lower return
volatility. A cooperative that prices to
members can absorb shocks through the surplus it would otherwise return, which
\citet{hesse2007} identify as the source of cooperative stability. For provisioning and
the intermediation margin it gives no reason to expect a difference. Minimum provisions
are fixed by a regulatory AA to H rating grid \citep{cmn2682}, independent of ownership,
so a difference on either is more likely to reflect the composition of the comparison
group. These expectations follow from the framework and are stated before the results;
the study was not pre-registered.

We hold no prior on the second question, which concerns method: which institutions
form the comparison group, how cumulative accounting flows are converted into rates, and
how a matched estimator is weighted.

% =============================================================================
\section{Methods}
% =============================================================================

\subsection{Data and classification}

We construct an institution-level quarterly panel from IF.data, the public statistical
repository of the Central Bank of Brazil. The summary financials, the segmentation file,
the asset composition file and the income statement are merged on institution code and
reporting date. Each source is deduplicated on the institution-quarter key before
merging, keeping the most complete record, and the raw panel runs from 2014Q1 to 2024Q4.
The window includes the pandemic years; quarter fixed effects absorb the common shock,
and the pre-2022 subperiod that contains them shows every surviving result
(Section~\ref{sec:robustness}). All tables and figures are the authors' calculations
from the IF.data files.

The panel is IF.data's prudential-conglomerate report, and every row carries its
reporting level. Of the banks, 87 report as consolidations and 57 individually, and some
switch within the window. All cooperatives report individually; each bank takes its modal
level, which is used below as a comparison-group parameter.

We classify institutions on the Central Bank's consolidation type, \textit{tipo de
consolidado banc\'ario} (\texttt{tcb}). The code \texttt{b3S} marks a singular credit
cooperative and \texttt{b3C} a central cooperative or confederation. Consolidation type
is a regulatory attribute and is stable over time; assigning each institution its modal
\texttt{tcb} makes cooperative status time-invariant by construction. Central
cooperatives are excluded, since they are wholesale entities that supply liquidity and
risk management to affiliated singulars and carry high capital for structural reasons. By
contrast, classification by name would place the sector's own banks, Banco Cooperativo
Sicredi and Banco Sicoob, both \texttt{b1} commercial banks, in the cooperative group.

\subsection{Comparison groups}
\label{sec:peers}

The cooperative side is always the singular retail cooperatives under the full
methodology. The comparison side is a reported parameter, defined on the regulator's
consolidation type and never on an outcome. The ladder has four rungs, and a fifth
comparison keeps the bank group but matches on macro-region as well as size
(Section~\ref{sec:region}).
\begin{itemize}\setlength{\itemsep}{0pt}
\item \textit{All non-cooperatives}: every full-methodology institution outside the
\texttt{b3} family.
\item \textit{Banks}, the primary group: \texttt{b1} commercial and multiple banks with a
commercial portfolio and \texttt{b2} multiple and investment banks.
\item \textit{Commercial}: \texttt{b1} banks only.
\item \textit{Structural}: \texttt{b1} banks with strictly positive deposits in every
quarter of their presence in the sample, defined ex ante.
\end{itemize}
Each rung can further be restricted to banks reporting individually. Of the 339
institutions in the broad group outside the bank rung, 158 have a median of zero credit
and zero funds raised (Figure~\ref{fig:broad}), which makes the credit and funding
differentials mechanical; we still report the broad group because its divergence from
banks is one of our results.

\begin{figure}[htbp]
\centering
\includegraphics[width=0.72\textwidth]{figures/fig14_broad_group.png}
\caption{Credit and funding ratios of full-methodology institutions}
\label{fig:broad}
\par\smallskip{\footnotesize\raggedright Institution medians over 2017 to 2024. The marker at the origin stands for 158 institutions. Values above 1.2 are drawn at 1.2.\par}
\end{figure}

\subsection{Outcomes and two measurement corrections}
\label{sec:measurement}

We study eight quarterly outcomes, defined in Table~\ref{tab:defs}, and add two stability
outcomes measured at the institution level over its full-methodology quarters: the
standard deviation of annualised quarterly return on assets and the z-score, the
distance-to-insolvency measure used by \citet{hesse2007}. Both require at least eight
quarters, which leaves 105 cooperatives and 137 banks, and the cooperatives dropped are
late entrants to the full methodology. Because the unit is the institution, these two
outcomes are estimated on one row per institution, with mean log assets as the size
control, no quarter effects, HC3 standard errors and a wild bootstrap over institutions.

Two features of IF.data affect how the quarterly outcomes are constructed. First, income
flows accumulate within the semester. Financial institutions in Brazil draw up balance
sheets semiannually, in June and December (Lei 4.595/1964), and the IF.data income fields
reset at each date. The median ratio of the fourth-quarter to the first-quarter value is
between 2.2 and 2.3 for every income field and close to one for balance-sheet stocks.
We de-cumulate within each semester and
annualise, which scales the return on assets coefficient by about 2.7 and leaves every
stock-based coefficient unchanged. Second, the intermediation result is reported net of
loan loss provisions, because the expense subtotal $(b)$ contains line $(b5)$,
\textit{Resultado de Provis\~ao para Cr\'editos de Dif\'icil Liquida\c{c}\~ao}, so $(c) =
(a) + (b)$ partly reproduces the provisioning result under another name. We compute the
margin gross of $(b5)$, and the cost-to-income denominator likewise.

The provisioning ratio is the stock of provisions constituted under Resolution 2.682/1999
over gross credit. The minimum provision is a step function of the AA to H rating
assigned to each operation, so the ratio reflects how the institution classifies its book
and is a weak proxy for realised delinquency.

Return on assets uses line $(j)$, net income after income tax, social contribution and
profit sharing. For a cooperative this is the surplus before its allocation to reserves
and to interest on member capital, which line $(k)$ reports separately. Because
cooperative acts are exempt from both taxes, Section~\ref{sec:results} re-estimates the
return on the pre-tax result, line $(g)$.

\begin{table}[htbp]
\centering
\small
\caption{Outcome definitions}
\label{tab:defs}
\begin{threeparttable}
\begin{tabular}{lp{10.2cm}}
\toprule
Outcome & Construction \\
\midrule
Basel capital ratio & \textit{Índice de Basileia} as reported by BACEN, in percentage
points \\
Leverage & Total liabilities / total assets \\
Return on assets & $4 \times$ line (j) / total assets \\
Intermediation margin & $4 \times [(c) - (b5)]\,/$ total assets \\
Cost-to-income & $[|(d3)| + |(d4)|] \,/\, [(c) - (b5) + |(d2)|]$ \\
Credit portfolio / assets & Credit portfolio / total assets \\
Provisioning ratio & Stock of provisions under Res.~2.682/1999 / gross credit \\
Funding ratio & \textit{Captações} (funds raised) / total assets \\
\midrule
Return volatility & S.d. of annualised quarterly return on assets over the institution's
full-methodology quarters \\
Z-score & [mean return on assets $+$ mean $(1 - \text{leverage})$] / s.d. of return on
assets \\
\bottomrule
\end{tabular}
\begin{tablenotes}[flushleft]\footnotesize
\item IF.data lines, de-cumulated within the semester: (a) intermediation revenue, (b) intermediation expense, carried as a negative, (c) $=$ (a) $+$ (b) intermediation result, (b5) loan-loss provision result within (b), (d2) fee
income, (d3) personnel expense, (d4) administrative expense, (j) net income.
\end{tablenotes}
\end{threeparttable}
\end{table}

\subsection{Specifications, matching and inference}
\label{sec:inference}

For each outcome we estimate the coefficient on a cooperative dummy under four
specifications: the unconditional difference in means, log total assets and quarter fixed
effects, coarsened exact matching \citep{iacus2012} with the same controls, and the same
controls on the region of common support in log assets. Cooperative status does not vary
within an institution, so it is collinear with any institution fixed effect. The
estimates are conditional correlations, and we do not present them as causal effects of
organisational form.

The coarsened-exact-matching estimand weights control units so that the
treated-to-control ratio is equal within every stratum \citep{iacus2012}; an unweighted
regression on the matched sample does not recover it, so we report both. We coarsen on
size quintile and control for region. Region-exact matching would leave too few
controls, since cooperatives concentrate in the South and banks in the Southeast: once
unequal weights are discounted, the effective control sample is about eighteen
institutions, against seventy-six under size coarsening. Region-exact estimates are
reported as a second parameter (Section~\ref{sec:region}).

Standard errors are clustered by institution, since cooperative status is assigned at the
institution level and observations are serially correlated \citep{cameron2015}.
Clustered inference is unreliable when few clusters are treated \citep{mackinnon2017},
as in the system-by-system estimates, so every coefficient also carries a restricted
wild cluster bootstrap, which re-draws the sign of each institution's residuals under
the null and compares the observed statistic with that distribution \citep{cameron2008}.
In simulations on our design with six treated clusters, the clustered test rejects a true
null 16 percent of the time at a nominal 5 percent level and 10 percent at a nominal 1
percent; the bootstrap rejects 10 percent and half of one percent, so its rejections at
one percent are conservative while those at five percent still over-reject.
System-level estimates report the number of treated clusters, and quantile
coefficients carry intervals from an institution bootstrap. Several outcomes are
right-skewed, so we report the conditional median alongside the conditional mean.

We classify each outcome by an explicit rule, reported in the last column of
Table~\ref{tab:main}. An outcome is \textit{null} unless it is significant at five
percent, with one sign, in all four specifications. A significant outcome is
\textit{stable} if it stays significant under every winsorisation rule, varies by less
than a factor of two across those rules, and has an average quantile coefficient between
the twentieth and eightieth percentiles of at least half the conditional mean. Otherwise
it is \textit{stable in sign}: the direction holds, but the tails or the shape of the
distribution drive the magnitude.

% =============================================================================
\section{Results}
\label{sec:results}
% =============================================================================

The analysis sample contains 130 singular cooperatives and 144 banks, 7,234
institution-quarters, from 2017Q1 to 2024Q4. Ninety percent of cooperative observations
fall inside the common support region of log assets (Figure~\ref{fig:overlap}), against
forty-eight percent of bank observations. About half the bank population therefore has no
cooperative counterpart at any size. Matching and trimming bring the standardised
difference in log assets from $-0.72$ to close to zero. Unconditional means and medians
are in the online appendix.

\begin{figure}[htbp]
\centering
\includegraphics[width=0.78\textwidth]{figures/fig8_overlap.png}
\caption{Log total assets, cooperatives and banks}
\label{fig:overlap}
\end{figure}

\subsection{The cooperative differential}

Table~\ref{tab:main} reports the coefficient on the cooperative dummy under the four
specifications, together with the conditional median counterpart of the second column and
the retention ratio $\rho$, the share of the raw differential that survives the
restriction to common support.

\begin{table}[htbp]
\centering
\begin{threeparttable}
\caption{Cooperative differential against banks}
\label{tab:main}
\input{tables/t2_main.tex}
\end{threeparttable}
\end{table}

Five of the eight quarterly outcomes keep their sign and significance across all four
specifications and every winsorisation rule, including no winsorisation. Their retention
ratios are above one, except for return on assets, which keeps most of its raw value; a
size-composition artefact would shrink as size is held more tightly, although the ratio
is not a formal test. By the rule of Section~\ref{sec:inference}, credit
intensity, return on assets and return volatility are stable. Capital, leverage and
cost-to-income are stable in sign, while provisioning, the margin and funding are null.
The capital coefficient should not be read at face value: under controls its conditional
mean is about seven times its conditional median, because the mean mostly reflects the
upper tail of the bank distribution, as we show below.

Cooperatives allocate a larger share of the balance sheet to credit, by nearly a fifth of
total assets under controls and by more at the conditional median. The differential moves
little across the four columns and changes by less than one percentage point across
winsorisation rules. Cooperatives also earn a higher return on assets, by about two
percentage points under controls and somewhat less at the conditional median, a large gap
by the standards of the banking literature. Their intermediation margin is
indistinguishable from that of banks, so the higher return appears to come from lending
volume and lower cost.

Part of the return gap is tax. Among quarters with a positive pre-tax result, the median
effective rate of income tax and social contribution is about a third for banks and
close to zero for cooperatives. Re-estimated on the pre-tax result, the controlled
differential is about three fifths of the after-tax figure and remains significant under
the bootstrap, and the conditional median keeps about five sixths of its value, but the
coefficient is marginal on the matched sample and not distinguishable from zero on
common support (Table~A11 of the online appendix). By the rule of
Section~\ref{sec:inference} the pre-tax differential is not established. Table~\ref{tab:main}
keeps the after-tax measure, the return that accumulates as capital, and we read the
profitability result as a difference after tax.

Cooperatives operate at lower cost, although the size of the difference is fragile. The
controlled coefficient roughly triples between tight winsorisation and none, because the
denominator approaches zero for banks that barely intermediate. We rely on the
conditional median, whose bootstrap interval excludes zero.

The capital result requires the most care. Unconditionally, the median cooperative is
better capitalised than the median bank, and Figure~\ref{fig:series} shows the
cooperative median above the bank median for most of the window. Yet the size-controlled
mean differential is large and negative in every column of Table~\ref{tab:main}.
Figure~\ref{fig:quantiles} reconciles the two: the cooperative coefficient is positive at
the bottom of the conditional capital distribution and negative at the top, and it
crosses zero near the fortieth percentile. Averaged over the twentieth to eightieth
percentiles, the differential is
about a third of the conditional mean (Table~\ref{tab:quantiles}). Unconditionally, 14
percent of bank-quarters exceed a Basel
ratio of 50 percent, against 1 percent of cooperative-quarters. At equal size, what mainly separates
the two groups is dispersion: residualised on size and quarter, the standard deviation of
the capital ratio among cooperatives is about a third of the bank figure (online
appendix), and the shaded bands of Figure~\ref{fig:series} show the same contrast in
every quarter.

\begin{table}[htbp]
\centering
\begin{threeparttable}
\caption{Quantile regression coefficients}
\label{tab:quantiles}
\input{tables/t7_quantiles.tex}
\end{threeparttable}
\end{table}

We examine two possible explanations of the bank upper tail. The first is reporting
level. The tail sits among the banks that report individually, whose ninetieth percentile
is far above that of the consolidations. The controlled gap against them is about four
times the gap against the consolidations (online appendix). The second is business model. Requiring a bank to be deposit-funded in every
quarter, the structural rung, roughly halves both the share of bank-quarters above 50
percent and the estimated mean gap, which remains significant (Figure~\ref{fig:ladder}).
The evidence is consistent with the bank upper tail being concentrated among small
standalone banks with little risk-weighted exposure. Against banks that lend and take
deposits, the capital difference consists of a modest gap in level and a large gap in
dispersion.

\begin{figure}[htbp]
\centering
\includegraphics[width=\textwidth]{figures/fig12_capital_quantiles.png}
\caption{The capital ratio across the distribution}
\label{fig:quantiles}
\par\smallskip{\footnotesize\raggedright Panel (a): quantile regression coefficients with a 95 percent institution-bootstrap band, and the conditional mean. Panel (b): Basel ratio before winsorisation, on a log scale; values below 5 percent are drawn at 5, and bar heights are percent of each group's quarters.\par}
\end{figure}

The same pattern holds for the other outcomes. Averaged over the twentieth to eightieth
percentiles, the differential is of similar size to the conditional mean for credit
intensity and return on assets, so for those two the mean describes the typical
institution. For leverage and cost-to-income it is little more than a third of the mean,
and both change sign inside the distribution. At equal size, cooperatives are the more
homogeneous group on every outcome, with residual standard deviations between a tenth and
a half of the bank figure. Leverage and capital are closely correlated among
cooperatives, and we treat the two as a single result.

Cooperative returns are also less volatile. At comparable size cooperatives show a lower
standard deviation of return on assets and a higher z-score, in all four designs and on
every national bank rung. Because the two measures share the standard deviation of return
on assets, we treat them as one result, with volatility as the primary measure. Together
with the return on assets result, this means that cooperatives earn higher after-tax and
less volatile returns on the same intermediation margin.

\subsection{Comparison-group dependence}

Table~\ref{tab:survival} states which differentials survive against which comparison
group, and Figure~\ref{fig:ladder} shows their size and precision on the same five
rungs: the broad group, banks, banks within macro-region, commercial banks and
structural banks. Return on assets, cost-to-income and return volatility survive on
every rung; credit intensity survives on the four national rungs and not within region
(Section~\ref{sec:region}). Dropping the sector's
two cooperative-owned banks changes every coefficient by 1.3 percent or less, apart from
the provisioning coefficient, which stays within 0.0002 of zero. Two results that appear
against the broad group disappear against banks.

\begin{table}[htbp]
\centering
\begin{threeparttable}
\caption{What survives against which comparison group}
\label{tab:survival}
\input{tables/t14_survival.tex}
\end{threeparttable}
\end{table}

\begin{figure}[htbp]
\centering
\includegraphics[width=\textwidth]{figures/fig13_ladder.png}
\caption{Cooperative coefficient by comparison group}
\label{fig:ladder}
\par\smallskip{\footnotesize\raggedright Controlled specification; the within-region rung also matches on macro-region. Intervals are 95 percent, clustered by institution (HC3 for the stability outcomes). Filled markers: $p < 0.05$. Shaded: banks, the primary group.\par}
\end{figure}

\begin{figure}[htbp]
\centering
\includegraphics[width=\textwidth]{figures/fig11_stable.png}
\caption{Medians and interquartile ranges, 2017 to 2024}
\label{fig:series}
\end{figure}

Against all non-cooperatives, cooperatives appear to provision less, significant at one
percent in all four specifications, while against banks the coefficient is close to zero.
The broad-group gap comes from the non-banking lenders, whose provisioning ratios are
nearly three times those of banks and cooperatives. At the
conditional median the sign reverses, with cooperatives provisioning more, and under the
tightest winsorisation rule the mean coefficient also turns positive and significant.
Supervisory concern attaches to capital, so we also split the sample into terciles of the
Basel ratio. No tercile shows cooperatives provisioning less: on the unconditional ratio
the differential is positive and significant only in the middle tercile, where most
cooperative quarters sit, and on the ratio residualised on size and quarter no tercile
differs. Because minimum provisions are a step function of the AA to H rating of
each operation, the result most likely describes how the two groups classify their
books.\footnote{Provisioning is observed for nearly all cooperative-quarters and for
about five in six bank-quarters. The missing bank-quarters are concentrated among the
low-credit institutions the ladder is designed to remove, and in 2023, when coverage falls
to about two thirds for banks.}

The intermediation margin behaves in the same way. It is large and highly significant
against the broad group, whose non-banking members report intermediation margins several
times those of banks and cooperatives relative to assets. Against banks it is small and
insignificant, and against structural banks it turns positive and significant. Its sign
depends on the comparison group, so we keep it as a null. The funding ratio is
significant in one of four specifications and reverses sign at the conditional median, so
we do not report it as a finding.

\subsection{Credit intensity within and across regional markets}
\label{sec:region}

Matched on size alone, cooperatives allocate a substantially larger share of assets to
credit than comparable banks, significant at one percent in every specification. Matched
also on macro-region, the differential is small and not distinguishable from zero under
any coarsening scheme. Cooperatives thus lend more than the national bank population,
while within regional markets the difference cannot be told apart from zero. However,
which estimate is the parameter of interest depends on an assumption we cannot test. If
cooperatives concentrate where they do for reasons internal to member-based finance,
region is a mediator and the size-only estimate is the parameter. If regional market
structure is heterogeneity unrelated to form, region is a confounder and credit intensity
is a null. We report the two as distinct parameters, the sector-wide and the
within-market differential, and label every claim by the comparison it rests on.

One caution applies. With an effective control sample of eighteen, the within-market
null may reflect low power as much as the parameter itself (Figure~\ref{fig:ladder}
shows the intervals). Return volatility survives the within-region comparison, while the z-score, which shares
its denominator, does not. We read the pair as one result whose within-market status is
uncertain.

\subsection{Network-level heterogeneity}

The online appendix splits the cooperative side by system. Sicredi and Sicoob account for
all but seven of the cooperatives, so the split describes those two networks and cannot
test network affiliation as such. Cooperatives in both networks hold less regulatory
capital than comparable banks, significant under the bootstrap in each. The six
independents rest on five treated clusters, and their capital estimate is positive and
cannot be separated from noise.

The split locates the pooled capital differential in the two large networks without
explaining it. A network backstop is one candidate explanation,
since a cooperative with centralised liquidity and consolidated risk management behind it
may run a thinner individual cushion. The networks also differ in the size, geography,
business mix and governance of their affiliates, and the split cannot separate the
backstop from these. Provisioning is insignificant in every system.

\subsection{Selection into the analysis sample}
\label{sec:selection}

The full-methodology restriction equalises the regulatory environment, but it also
selects the sample. The cooperatives we observe are the largest in the country, about a
tenth of all cooperatives recorded in the panel. Fifty-one singular cooperatives are
observed under the simplified methodology before they adopt the full methodology during
the window.\footnote{Five cooperatives with simplified quarters only after their first
full quarter are excluded from this comparison. With one adopter that also reverts, they
are the six reverters of Section~\ref{sec:robustness}; including them changes no sign.} We
compare those future adopters with cooperatives that never adopt, on pre-adoption
quarters only and conditioning on size and quarter. The two groups show no difference in
leverage, but adopters were less credit-intensive ($-0.068$, $p=0.002$) and marginally
less profitable, which runs against our findings. They were also marginally leaner, which
runs with them. As a result, selection likely biases the credit and profitability
estimates towards zero and, if anything, the efficiency estimate marginally away from it.

\subsection{Sensitivity checks}
\label{sec:robustness}

The online appendix reports the conditional medians and the specification without the
eleven S1 and S2 banks, which moves no significant coefficient by more than two percent.
In addition, it reports a split at the January 2022 replacement of the 2013 capital
rules. Every surviving result keeps its sign and significance in both subperiods, and
provisioning and the margin remain null in both. Magnitudes are roughly a fifth smaller
after 2022. Four winsorisation rules, re-estimation on the rows where provisioning is
observed, and dropping the six cooperatives that revert to the simplified methodology
leave every surviving sign unchanged. The weighted and unweighted matched-sample
estimates differ by less than the sampling error. All results
come from one script run on the public IF.data files, with every comparison group,
coarsening, winsorisation rule and check as a parameter, and the tables are generated
from its output.

% =============================================================================
\section{Discussion}
% =============================================================================

The surviving results suggest a focused retail lender. A cooperative in the
full-methodology population channels a high share of its assets into member credit and
runs that operation at low cost. It earns a solid after-tax return on assets on
intermediation spreads similar to those of banks, and it sits in a narrow band of the capital
distribution. This profile is consistent with a member-service orientation in which
capital accumulation is a means to member credit \citep{fonteyne2007}.

The capital results also qualify the property-rights argument. \citet{cook1995}
predicts weak incentives to build capital, yet the median cooperative is better
capitalised than the median bank and holds more capital at the thin end of the
distribution. The two fit together once the source of capital is considered. The
property-rights problems bear on capital raised from members, which a cooperative cannot
expand at will; the reserves it retains are indivisible and accumulate across member
generations \citep{fonteyne2007}. What the argument rules out is the upper tail, since a
cooperative has no reason to hold capital far above what its lending requires and no
outside investor to supply it. The narrow band is the profile that endowment capital
predicts, and the mean gap is the absence of that tail.

Our results agree with \citet{hesse2007} on two of the three bank-level differences
they report, lower capitalisation at comparable size and lower return volatility. They
differ on the third.
In their sample cooperative stability came with lower profitability, whereas in this
population cooperatives are more profitable after tax and less volatile, on margins
indistinguishable from those of banks. Part of the profitability gap is the tax
exemption of cooperative acts, and before tax the gap is smaller and not established,
which narrows the disagreement. The rest may reflect our population, which is restricted
to the largest cooperatives and to a common prudential methodology.

For supervision, the relevant combination is capital and provisioning taken together. At
the thin end of the capital distribution, where prudential concern would attach,
cooperatives hold more capital than comparable banks and provision no less, while the
finding that they provision more holds only in the middle of the capital distribution.
The data also show that cooperatives lack the highly capitalised upper tail that banks
exhibit. Several differentials that look large at the mean are differentials in
dispersion, so we give the conditional median and the quantile profile as much weight as
the conditional mean.

Three analyst choices determined whether headline results existed. The first is the
comparison group, which as ``all non-cooperatives'' imports some 160 institutions that
neither lend nor raise funds. The second and third are reading IF.data income fields as
quarterly flows and estimating the matched comparison by unweighted OLS on the matched
sample. Each is a natural reading of the source material. Some of the disagreements in
the literature may reflect different comparison groups.

\subsection{Limitations}

Six limitations bound the claims. Because cooperative status is time-invariant, the
estimates are conditional correlations. The full-methodology restriction selects the
largest cooperatives, and the results generalise only to that subpopulation. The
provisioning ratio is an accounting variable and does not measure delinquency. The
stability outcomes are institution-level statistics on a subsample, estimated without
quarter effects. Most systems are too thinly represented for separate inference, and the
network backstop hypothesis needs consolidated network data that IF.data does not carry.
The mediator question in Section~\ref{sec:region} also needs a market measure finer than
macro-region, and IF.data's credit-modality modules, which carry rural credit, are the
natural next step. Finally, a within-institution design around the adoption of the full
methodology does not identify a regime effect: adoption follows growth, which moves every
balance-sheet ratio continuously, and pre-adoption trends run through the adoption date
without a break.

% =============================================================================
\section{Contributions and Implications}
% =============================================================================

This paper characterises cooperative behaviour under full Basel-style regulation in a
large emerging market, on the full population of full-methodology institutions. The
substantive contribution is this profile, together with two nulls that correct claims the
same data appear to support under a looser comparison group. The methodological
contribution is a protocol that defines peers on the regulator's own categories, reports
a ladder of comparison groups, and separates sector-wide from within-market differences.
It also states how cumulative flows were converted to rates, reports the matched
estimand and prints treated cluster counts beside system-level $p$-values; each step is
inexpensive, and each changed at least one conclusion.

Our results bear on two margins of proportionality beyond size \citep{cmn4553}. Because
the evidence is descriptive, we frame both as questions. The first is organisational
architecture, since the capital differential is a property of the two large networks.
Brazil already treats the cooperative sector as a unit for deposit insurance through the
FGCoop. If the backstop hypothesis is right, a proportional framework might also treat
an integrated network and its backstop as the unit for parts of the capital and
liquidity assessment, as \citet{coelho2019} report for the French cooperative groups. If
the differential
reflects the affiliates' size, geography or business mix, that suggestion does not
follow. The second is the calibration of the benchmark itself. A benchmark set against
the average full-methodology institution may describe the cooperative configuration
poorly, because that average comes from a distribution whose upper region cooperatives do
not inhabit. Section~\ref{sec:theory} indicates why a lighter requirement does not follow
automatically. The representation rationale weakens for a member-owned, network-monitored
cooperative, but the systemic externality and the guarantee remain, and the FGCoop's
flat-rate mutual funding places the cost of a thin cushion on the rest of the sector, so
thinner capital may call for closer supervisory attention as well as for relief.

For the cooperative movement, the difference between the two forms persists under a
common prudential methodology, in a narrower form than a loose comparison suggests. Our
evidence does not support a case resting on cooperatives provisioning more conservatively
than banks. A case resting on cost at the median and on return stability holds against
every definition of a bank we tried; one resting on the return on assets holds after
tax, and before tax the gap is smaller and not established. One resting on credit
intensity holds against the national bank population but has not been shown within
regional markets; the territorial claim is the one these data cannot yet settle, and the
rural-credit modules of IF.data are the next step.

\subsection*{Declaration}

The authors work at the Organization of Brazilian Cooperatives (OCB System), which
represents the cooperative sector in Brazil. The study uses public data and a public
reproduction pipeline, and the views expressed are the authors' own.

% =============================================================================
\section{Key References}
% =============================================================================

\renewcommand{\refname}{}
\vspace{-2.5em}
\bibliographystyle{unsrtnat}
\bibliography{references}

\end{document}
